\documentclass[12pt]{article}
\usepackage[utf8]{inputenc}
\usepackage[T1]{fontenc}
\usepackage{hyperref}
\title{Tswapong Language Resources for K-12 Curriculum Development}
\author{Compiled by Claude Sonnet 5 (claude-sonnet-5) for Andr\'as Kornai}
\date{}
\begin{document}
\maketitle

\section{Introduction}
Tswapong (Setswapong; ISO 639-3 two, Glottolog tswa1237) is a Sotho-Tswana (Bantu) language of eastern Botswana, around the Tswapong Hills. The Endangered Languages Project lists it as threatened, citing about 2,000 speakers from a 1997 survey (Andersson and Janson). Joshua Project gives a much larger figure for the people group, about 9,500 to 9,600. I did not reconcile the two. I found no grammar, dictionary or school materials.

\section{Primary Educational Resources}
None found.

\section{Linguistic and Reference Materials}
\begin{itemize}
\item Endangered Languages Project entry: \url{https://www.endangeredlanguages.com/lang/3813}
\item Verbix endangered-languages page for Tswapong (contents not checked): \url{https://docs.verbix.com/EndangeredLanguages/Tswapong}
\end{itemize}

\section{Religious and Community Resources}
\begin{itemize}
\item Botswana Government news report (June 2026): a Setswapong Bible translation project, 2026 to 2028, run by the Batswapong Cultural and Development Society with the Botswana Bible Society. The report also mentions an advisory committee working on a harmonised written form, and a broadcasting licence for a community radio station, Tswapong FM, at Lerala. \url{https://dailynews.gov.bw/news-detail/92023}
\item Joshua Project people profile: \url{https://joshuaproject.net/people_groups/15602}
\item Joshua Project language page: \url{https://joshuaproject.net/languages/two}
\end{itemize}

\section{Digital and Audio Resources}
Tswapong FM (planned community radio, per the news report above). No recordings found in search results.

\section{Curriculum Development Context}
The reference description says the language is not known to be taught in schools. The standardisation work and the Bible project, which are new, could supply a written norm and text. Setswana is the likely school language. A curriculum would need new materials developed with the Batswapong Cultural and Development Society.
\end{document}